\documentclass[12pt,a4paper,italian]{article}


\usepackage[italian]{babel}
\usepackage[latin1]{inputenc}
\usepackage{amsmath}
\usepackage{amsfonts}
\usepackage{amssymb}
\usepackage{color}
\usepackage{xcolor}
\usepackage{hyperref}
\usepackage[all]{hypcap}
\usepackage{ifthen}
\usepackage{wrapfig}
%\author{Piero Bizzotto}
\usepackage[top=2cm,bottom=5cm,left=80pt,right=80pt]{geometry}
\usepackage{graphicx}
\DeclareGraphicsExtensions{.jpg,.png}

\newcommand{\ajax}{AJAXDRAW}
\newcommand{\sito}{\href{http://ajaxdraw.sourceforge.net}{http://ajaxdraw.sourceforge.net}}

\setlength{\parindent}{0pt} %settato indentazione di default 
\setlength{\headheight}{3cm} %settato grandezza header...in altre parole, quanto distanzio il doc dall'intestazione

\usepackage{fancyhdr} %pacchetto per le intestazioni
\pagestyle{fancy} %uso del pacchetto


\fancyhead{} %annulla head di default
\fancyfoot{} %annulla foot di default


\usepackage{lastpage} %setto pg di pgtot a rfoot
     \rfoot{pagina \thepage\ di \pageref{LastPage}}


\lfoot{Versione: \insertversion} %setto versione doc a lfoot
\renewcommand{\footrulewidth}{0.5pt} %ridefinisco il valore della riga di intestazione
\renewcommand{\headrulewidth}{0.5pt} %ridefinisco il valore della riga di pie' di pagina

\newcommand{\insertversion}{0.0} %definisco il nuovo comando per inserire la versione


\lhead{  \begin{Huge} \ajax \end{Huge} \\  %intestazione di sinistra
					%\begin{Large}	Software per il Disegno Grafico\\ in Tecnologie Web \end{Large}  
					\begin{normalsize}\sito \end{normalsize}
			%\\ versione documento: \insertversion\ del \today} %setto l'intestazione sx
		}
\rhead{ %includo logo nell'intestazione dx
	 	\includegraphics[scale=0.5]{../logo/logo}  
}


%CREAZIONE ELENCHI NUMERATI PERSONALIZZATI
\newcounter{Lcount}
\newcounter{Rcount}
\setcounter{Lcount}{0}
\setcounter{Rcount}{0}

\newenvironment{elenconumerato}[2][ ]
{
  \begin{list}{#1\arabic{Lcount}.}
    {
	\setcounter{Rcount}{\value{Lcount}}
	\setcounter{Lcount}{0} 
	\usecounter{Lcount} 
\addtolength{\leftmargin}{#2pt}
	}
}
{
  \end{list}
 \setcounter{Lcount}{\value{Rcount}}
}

%CREAZIONE ELENCHI PUNTATI
\newenvironment{elencopuntato}[1][]
{
\begin{list}{\textbullet} %\itemindent=#1pt
	{
	\addtolength{\leftmargin}{#1pt}
	}
} 
{
\end{list}
}


\newenvironment{elencodescrittivo}[1][]{\begin{description} \setlength{\itemindent}{#1pt} \addtolength{\leftmargin}{#1pt}} {\end{description}}

\newcommand{\TITOLODOC}{Titolo}

%footer centrale
\cfoot{ \TITOLODOC \\  E-mail:    \href{ mailto:webshape.contact@gmail.com}{ webshape.contact@gmail.com}  }

%INSERIMENTO IMMAGINI
\newcommand{\imagerealsize}[1]{\vspace{20pt} \includegraphics{#1} }
\newcommand{\imageadapted}[1]{\vspace{20pt} \includegraphics[width=1\textwidth]{#1} }

\newcommand{\glosspath}{.\glossario}
\newcommand{\gloss}[1]{\hyperref{\glosspath~\glossario.pdf}{}{#1}{#1}}

\hypersetup{
    %bookmarks=true,         % show bookmarks bar?
    %unicode=false,          % non-Latin characters in Acrobat’s bookmarks
	%pdftoolbar=true,        % show Acrobat’s toolbar?
	%pdfmenubar=true,        % show Acrobat’s menu?
    %pdffitwindow=true,      % page fit to window when opened
    %pdftitle={My title},    % title
    %pdfauthor={Author},     % author
    %pdfsubject={Subject},   % subject of the document
    %pdfnewwindow=true,      % links in new window
    %pdfkeywords={keywords}, % list of keywords
    colorlinks=true,         % false: boxed links; true: colored links
    linkcolor=black,           % color of internal links
    %citecolor=green,        % color of links to bibliography
    %filecolor=magenta,      % color of file links
    urlcolor=teal    % color of external links
%	linktocpage=false;
}


%COLORAZIONE TESTO
\newcommand{\blue}[1]{{\color {blue} #1}} 
\newcommand{\red}[1]{{\color {red} #1}}
\newcommand{\green}[1]{{\color {green} #1}}
\newcommand{\sezione}[1]{\leftskip=0pt \section{#1} \leftskip=18pt}
\newcommand{\subsezione}[1]{\leftskip=18pt \subsection{#1} \leftskip=36pt}
\newcommand{\subsubsezione}[1]{\leftskip=36pt \subsubsection{#1} \leftskip=54pt}
\newcommand{\subsubsecindent}{54}
\newcommand{\subsecindent}{36}
\newcommand{\secindent}{18}
\newcommand{\normindent}{8}
\newcommand{\code}[1]{{\bfseries \texttt{#1}}}
\newcommand{\paragrafo}[1]{\leftskip=36pt \paragraph{#1} \leftskip=54pt}
\newcommand{\subparagrafo}[1]{\leftskip=54pt \subparagraph{#1} \leftskip=72pt} %BASE!!!
\usepackage{multirow}
\begin{document}

\renewcommand{\insertversion}{0.0} %INSERIRE LA VERSIONE QUI DENTRO STILE x.x.xx
\renewcommand{\TITOLODOC}{Analisi dei Requisiti} %INSERIRE IL TITOLO DEL DOCUMENTO DA FAR COMPARIRE A PIE PAGINA

\begin{titlepage}
\begin{center}
	\begin{Large}	\today \end{Large}
\end{center}

\vspace{20pt}

\begin{center}
	\begin{Huge}
				\textbf{AJAXDRAW}
	\end{Huge}
\end{center}			

\begin{center}
	\begin{large}
				\textbf{Software per il Disegno Grafico\\ in Tecnologie Web}
	\end{large}
\end{center}			

\vspace{20pt}

\begin{center}
\includegraphics[width=150pt]{../logo/logo}
\end{center}

\vspace{160pt}
\begin{center} %INSERIRE ALL'INTERNO IL TITOLO DOCUMENTO CHE COMPARIRA NELLA PAGINA INIZIALE				
	\begin{Huge}
				\textbf{\TITOLODOC}
	\end{Huge}
			\\
\end{center}
\vspace{220pt}
\begin{center}
Versione: \insertversion
\end{center}
\end{titlepage}

\newpage


\begin{center} %INSERIRE ALL'INTERNO IL TITOLO DOCUMENTO CHE COMPARIRA NELLA PAGINA INIZIALE
	\begin{Huge}	
				\textbf{\TITOLODOC}
			\\
	\end{Huge}
\end{center}
\parindent=18pt %settato indentazione di default 
\section*{\LARGE Sommario:} %SEZIONE SOMMARIO
Questo documento si prefigge di presentare lo studio effettuato da WebShape riguardo al prodotto software relativo al capitolato d'appalto denominato C04 (\blue{\ajax}, Software per il Disegno Grafico in Tecnologie Web).\\
Tale studio \`e mirato alla comprensione dei bisogni espressi nel capitolato e alla loro formalizzazione e classificazione in requisiti informatici. In particolare quelli funzionali individuati nel documento verranno espressi anche mediante casi d'uso, sia narrativi che grafici. Questo documento avr\`a valore contrattuale.

\section*{\LARGE Stato del documento:}
	Formale Esterno
\hangindent=0pt

\section*{\LARGE Redazione:}
	\begin{table}[!h]
		\begin{center}
			\begin{tabular}
				{|c|c|}
				\hline
				%%%%%%%%%%%%%%INTESTAZIONE COLONNE%%%%%%%%%%%%%%%%%%%%%%%%%%%%%%%%%%%%%%%%%%%%%%%%%%%%%%%%%%%%%%
				\multicolumn{2}{|c|}{ \textbf{Redazione} } \\
				\hline
				\textbf{Fase} & \textbf{Redattori} \\
				%%%%%%%%%%%%%%FINE INTESTAZIONE COLONNE%%%%%%%%%%%%%%%%%%%%%%%%%%%%%%%%%%%%%%%%%%%%%%%%%%%%%%%%%%%%%%
				\hline
				%%%%%%%%%%% PARTE DA MODIFICARE %%%%%%%%%%%%%%%%%%%%%%%%%%%%%%%%%%%%%%%%%%%%%%%%%%%%%%%%%%%%%%%%%		
				\multirow{2}{*}{Pre-RR} & Bizzotto Piero\\
										& Dissegna Stefano\\
				\hline
				\multirow{2}{*}{RR-RPP} & \\
										& \\
				\hline
				%%%%%%%%%%% FINE PARTE DA MODIFICARE %%%%%%%%%%%%%%%%%%%%%%%%%%%%%%%%%%%%%%%%%%%%%%%%%%%%%%%%%%%
			\end{tabular}
			\caption{Lista Redattori} %INSERIRE DIDASCALIA - SE NECESSARIA - 
			\label{tabredazione}
		\end{center}
	\end{table}	
\section*{\LARGE Approvazione:}
%	\begin{elencopuntato}[\normindent]
%		\item 
%		\item 
%	\end{elencopuntato}
\begin{table}[!h]
	\begin{center}
		\begin{tabular}
			{|c|c|}
			\hline
			%%%%%%%%%%%%%%INTESTAZIONE COLONNE%%%%%%%%%%%%%%%%%%%%%%%%%%%%%%%%%%%%%%%%%%%%%%%%%%%%%%%%%%%%%%
			\multicolumn{2}{|c|}{ \textbf{Approvazione} } \\
			\hline
			\textbf{Fase} & \textbf{Approvatori} \\
			%%%%%%%%%%%%%%FINE INTESTAZIONE COLONNE%%%%%%%%%%%%%%%%%%%%%%%%%%%%%%%%%%%%%%%%%%%%%%%%%%%%%%%%%%%%%%
			\hline
			%%%%%%%%%%% PARTE DA MODIFICARE %%%%%%%%%%%%%%%%%%%%%%%%%%%%%%%%%%%%%%%%%%%%%%%%%%%%%%%%%%%%%%%%%		
			\multirow{2}{*}{Pre-RR} &  \\
									&  \\
			\hline
			\multirow{2}{*}{RR-RPP} & \\
									& \\
			\hline
			%%%%%%%%%%% FINE PARTE DA MODIFICARE %%%%%%%%%%%%%%%%%%%%%%%%%%%%%%%%%%%%%%%%%%%%%%%%%%%%%%%%%%%
		\end{tabular}
		\caption{Lista Approvatori} %INSERIRE DIDASCALIA - SE NECESSARIA - 
		\label{tabapprovazione}
	\end{center}
\end{table}
\textbf{}

\section*{\LARGE Lista di Distribuzione:}

	\begin{elenconumerato}{\normindent}
		\item WebShape \footnote{Il termine WebShape designa una collettivit\`a di individui come da organigramma contenuto nel piano di progetto fornito in allegato al presente documento}
		\item I committenti Conte Renato e Vardanega Tullio in rappresentanza \\  dell'azienda proponente Zucchetti SPA
	\end{elenconumerato}

\newpage

\section*{\LARGE Registro delle Modifiche:}

\begin{center}
	\begin{table}[h]
		  \begin{tabular*}
			{1\textwidth}%
					 {@{\extracolsep{\fill}}|p{0.1\textwidth}|p{0.55\textwidth}|p{0.25\textwidth}|}
		 \hline
%%%%%%%%%%%%%%INTESTAZIONE COLONNE%%%%%%%%%%%%%%%%%%%%%%%%%%%%%%%%%%%%%%%%%%%%%%%%%%%%%%%%%%%%%%
			\textbf{Versione}  & \textbf{Descrizione} & \textbf{Autore} \\
%%%%%%%%%%%%%%FINE INTESTAZIONE COLONNE%%%%%%%%%%%%%%%%%%%%%%%%%%%%%%%%%%%%%%%%%%%%%%%%%%%%%%%%%%%%%%
		 \hline
%%%%%%%%%%% PARTE DA MODIFICARE %%%%%%%%%%%%%%%%%%%%%%%%%%%%%%%%%%%%%%%%%%%%%%%%%%%%%%%%%%%%%%%%
		
                \hline
                0.1 & 30$\slash$11$\slash$2008 Casi d'uso relativi ad avvio e interazione con l'applicazione. Aggiunti requisiti. & Dissegna Stefano \\
		\hline	
    	 	0.0 & 		 26$\slash$11$\slash$2008 Strutturazione del documento. & Bizzotto Piero \\

		\hline %%FINE RIGA
%%%%%%%%%%% FINE PARTE DA MODIFICARE %%%%%%%%%%%%%%%%%%%%%%%%%%%%%%%%%%%%%%%%%%%%%%%%%%%%%%%%%%%
		\end{tabular*}
	\caption{Registro delle modifiche} %INSERIRE DIDASCALIA - SE NECESSARIA - 
	\label{tab:modifiche}
	\end{table}
\end{center}

\newpage
\thispagestyle{fancy}
\tableofcontents
\thispagestyle{fancy}
\newpage
\parskip=-5pt

\sezione{Introduzione}

\subsezione{Scopo del documento}
Il presente documento \`e indirizzato a fornire una descrizione in grado di identificare il prodotto software AJAXDRAW, soggetto a gara d'appalto a cui WebShape intende concorrere.\\
Il documento elenca pertanto i requisiti, impliciti, espliciti, funzionali e non funzionali, individuati per il prodotto di cui sopra.
%Si intende anche fornire una base per definire gli accordi con il cliente.

\subsezione{Scopo del prodotto}
AJAXDRAW \`e inteso principalmente come "Proof of Concept", ossia un'incompleta realizzazione di uno strumento grafico completamente basato sul web, con lo scopo di dimostrarne la fattibilit\`a o la fondatezza di alcuni principi o concetti costituenti. In particolare permetter\`a a chiunque di poter disegnare, tramite dei tools appositamente forniti, delle immagini sulla falsariga del software open-source per il disegno vettoriale Inkscape.

\subsezione{Riferimenti normativi}

Il presente documento \`e redatto in accordo con le norme interne di WebShape, raccolte nel documento NormeDiProgetto.pdf, consultabile per conoscenza dal repository pubblico al quale WebShape si appoggia per i suoi progetti.

\sezione{Descrizione generale}

\subsezione{Contesto d'uso del prodotto}

\subsubsezione{Processi produttivi e modalit\`a d'uso}
Il sistema sar\`a costituito da una pagina web eseguibile con uno qualsiasi dei pi\`u diffusi browsers attualmente esistenti.

\subsubsezione{Piattaforme d'esecuzione,\\ interfacciamento con l'ambiente di installazione e uso}
Il prodotto \`e destinato all'uso da parte di qualunque persona abbia installato nel proprio computer uno dei principali internet browsers, quindi \`e possibile garantire il funzionamento su ogni ambiente di esecuzione. (Si rimanda alla lista dei requisiti in sezione \ref{listarequisiti}).

\subsezione{Funzioni del prodotto}
AJAXDRAW permetter\`a di avere un software di disegno grafico in grado di poter elaborare figure vettoriali e complesse direttamente sul proprio browser web. Le caratteristiche di tale applicazione saranno facilmente utilizzabili dall'utente, che in particolare potr\`a scegliere, tramite pulsanti posizionati su una toolbar, quali azioni effettuare, come disegno di primitive, inserimento di testo, spostamento degli oggetti.\\
WebShape si riserva di decidere in un secondo momento se implementare funzionalit\`a aggiuntive pi\`u complesse.\\
L'applicativo permetter\`a l'esportazione del disegno creato nel formato SVG, diventato una raccomandazione (standard) del World Wide Web Consortium.

\subsezione{Caratteristiche degli utenti}
\label{definizione_utente}
Gli utenti previsti per AJAXDRAW sono normali utilizzatori di Internet, quindi persone dalle varie capacit\`a tecniche, dalle pi\`u basse alle pi\`u alte. Si assume che l'utente finale conosca le funzioni basilari di un software di grafica, quindi si danno per noti concetti di disegno di primitive, inserimento di testo, selezione di oggetti, in quanto contenuti al giorno d\'oggi nella stragrande maggioranza dei software pi\`u usati.

\subsezione{Vincoli generali}
Il software che WebShape si impegna a sviluppare sar\`a soggetto a licenza \blue{GPL}.

\subsezione{Dipendenze}
Si assume che i sistemi in cui verr\`a eseguito AJAXDRAW siano dotati di uno tra i browsers Mozilla Firefox, Google Chrome, Apple Safari, Opera e Internet Explorer nelle loro pi\`u recenti incarnazioni.

\sezione{Glossario}
Come previsto dalle norme interne di WebShape, tutti i documenti fanno riferimento ad un unico glossario (Glossario.), fornito in allegato al presente documento.

\sezione{Casi d'uso}
Vengono riportati di seguito i casi d'uso, sia grafici che narrativi, utilizzati durante l'analisi dei requisiti effettuata da WebShape. Lo scopo dei casi d'uso \`e quello di permettere una chiara e semplice comprensione dei bisogni del cliente percepiti dall'azienda. Per una lista esaustiva dei requisiti individuati si rimanda comunque alla lista dei requisiti di sezione \ref{listarequisiti}, cui faranno riferimento anche gli stessi casi d'uso narrativi.

\subsezione{Avvio}
Il seguente grafico descrive come ci si attende l'utente avvii l'applicazione.
\begin{figure}
\centering
%\includegraphics{}
\caption{avvio dell'applicazione}
\end{figure}

Segue ora una descrizione testuale del caso d'uso presentato dal grafico.

\subsubsezione{UC0 Caso d'uso: richiedi pagina web}
\paragraph{Attori coinvolti} Utente, browser, server HTTP
\paragraph{Scopo e descrizione sintetica}
L'utente avvia l'applicazione visualizzandone la relativa pagina web attraverso il browser, che richiede la relativa pagina al server HTTP.
\paragraph{Flusso di eventi}
\begin{enumerate}
\item L'utente inserisce l'URL relativo all'applicazione nel browser.
\item Il browser richiede la pagina al server HTTP
\item Il server HTTP invia la pagina.
\item Il browser visualizza la pagina.
\end{enumerate}
\paragraph{Flusso alternativo}
Se viene rilevato un browser non supportato, l'utente viene avvisato. L'utente pu\`o decidere di continuare lo stesso senza nessuna garanzia di corretto funzionamento.
\paragraph{Precondizioni} Il browser \'e avviato.
\paragraph{Postcondizioni} L'applicazione \'e avviata.

\subsezione{Interazione con l'applicazione}
Il seguente grafico descrive le interazioni tra l'utente finale e l'applicazione da un punto di vista ad alto livello. Le interazioni principali verranno poi espanse.
\begin{figure}[!ht]
\centering
\vspace{20pt} 
%\includegraphics[width=1\textwidth]{UC0}
\caption{funzionalit\`a dell'applicazione}
\end{figure}

\paragraph{Segue ora una descrizione testuale dei casi d'uso presentati dal grafico.}

\subsubsezione{UC1 Caso d'uso: pulisci \underline{canvas}}
\paragraph{Attori coinvolti} Utente
\paragraph{Scopo e descrizione sintetica}
L'utente pu\`o cancellare totalmente quanto finora disegnato.
\paragraph{Flusso di eventi}
\begin{elenconumerato}[\textbf{}]{\subsubsecindent}
\item L'utente seleziona il comando per pulire il canvas. 
\item Il canvas viene completamente cancellato.
\end{elenconumerato}
\paragraph{Precondizioni} Il canvas \'e visibile.
\paragraph{Postcondizioni} Il canvas \'e vuoto.

\subsubsezione{UC2 Caso d'uso: visualizza manuale}
\paragraph{Attori coinvolti} Utente
\paragraph{Scopo e descrizione sintetica}
L'utente pu\`o visualizzare il manuale dell'applicazione in linea.
\paragraph{Flusso di eventi}
\begin{elenconumerato}[\textbf{}]{\subsubsecindent}
\item L'utente seleziona il comando per visualizzare il manuale. 
\item Il manuale viene visualizzato
\end{elenconumerato}
\paragraph{Precondizioni}
\paragraph{Postcondizioni} Il manuale utente \'e visualizzato

\subsubsezione{UC3 Caso d'uso: richiedi manuale}
\paragraph{Attori coinvolti} Server HTTP
\paragraph{Scopo e descrizione sintetica}
La/le pagina/e web contenenti il manuale utente sono gestite dal server HTTP.
\paragraph{Flusso di eventi}
\begin{elenconumerato}[\textbf{}]{\subsubsecindent}
\item Il sistema richiede la pagina web contenente il manuale utente al server HTTP.
\item Il server HTTP ritorna la pagina richiesta
\end{elenconumerato}
\paragraph{Flusso alternativo}
Se il server HTTP non \'e disponibile o vi \'e un errore di comunicazione, l'utente viene avvisato.
\paragraph{Precondizioni} L'utente ha richiesto al sistema la visualizzazione del manuale o di parte di esso.
\paragraph{Postcondizioni} Il sistema dispone della pagina web richiesta.

\subsubsezione{UC4 Caso d'uso: sposta visualizzazione}
\paragraph{Attori coinvolti} Utente
\paragraph{Scopo e descrizione sintetica}
Quando il canvas \'e troppo grande per essere completamente visualizzato, l'utente pu\`o spostare la visualizzazione per vedere le diverse parti del canvas.
\paragraph{Flusso di eventi}
\begin{elenconumerato}[\textbf{}]{\subsubsecindent}
\item L'utente seleziona la parte del canvas da visualizzare
\item Il sistema visualizza la parte di canvas selezionata al centro dello schermo.
\end{elenconumerato}
\paragraph{Precondizioni} Il canvas non \'e completamente visualizzato a schermo.
\paragraph{Postcondizioni} L'utente vede su schermo la parte di canvas desiderata.

\subsubsezione{UC5 Caso d'uso: aumenta zoom}
\paragraph{Attori coinvolti} Utente
\paragraph{Scopo e descrizione sintetica} 
Il sistema permette all'utente di ingrandire il canvas ed il suo contenuto mantendendo le proporzioni.
\paragraph{Flusso di eventi}
\begin{elenconumerato}[\textbf{}]{\subsubsecindent}
\item L'utente seleziona, esplicitamente o implicitamente, la parte di canvas da mantenere al centro della visualizzazione.
\item L'utente seleziona di quanto aumentare lo zoom.
\item Il sistema ingrandisce il canvas di quanto richiesto
\item Il sistema accentra la visualizzazione dove richiesto dall'utente
\end{elenconumerato}
\paragraph{Precondizioni} Il canvas \'e visualizzato
\paragraph{Postcondizioni} Il canvas \'e ingrandito

\subsubsezione{UC6 Caso d'uso: diminuisci zoom}
\paragraph{Attori coinvolti} Utente
\paragraph{Scopo e descrizione sintetica} 
Il sistema permette all'utente di rimpicciolire il canvas ed il suo contenuto mantendendo le proporzioni.
\paragraph{Flusso di eventi}
\begin{elenconumerato}[\textbf{}]{\subsubsecindent}
\item L'utente seleziona, esplicitamente o implicitamente, la parte di canvas da mantenere al centro della visualizzazione.
\item L'utente seleziona di quanto diminuire lo zoom.
\item Il sistema rimpicciolisce il canvas di quanto richiesto.
\item Il sistema accentra la visualizzazione dove richiesto dall'utente.
\end{elenconumerato}
\paragraph{Precondizioni} Il canvas \'e visualizzato.
\paragraph{Postcondizioni} Il canvas \'e rimpicciolito.

\subsubsezione{UC7 Caso d'uso: esporta in SVG}
\paragraph{Attori coinvolti} Utente
\paragraph{Scopo e descrizione sintetica} 
Il sistema permette all'utente di esportare in formato SVG il contenuto del canvas.
\paragraph{Flusso di eventi}
\begin{elenconumerato}[\textbf{}]{\subsubsecindent}
\item L'utente seleziona il comando per esportare il canvas in formato SVG.
\item Il sistema genera un file SVG il cui contenuto corrisponde a quanto correntemente visualizzato nel canvas.
\item Il sistema rende disponibile all'utente il file generato.
\end{elenconumerato}
\paragraph{Precondizioni} Il canvas \'e visualizzato.
\paragraph{Postcondizioni} Il file SVG \'e disponibile all'utente e contiene quanto visualizzato nel canvas.

\subsubsezione{UC8 Caso d'uso: carica file SVG}
\paragraph{Attori coinvolti} Utente
\paragraph{Scopo e descrizione sintetica} 
Il sistema permette all'utente di caricare file in formato SVG nel canvas.
\paragraph{Flusso di eventi}
\begin{enumerate}
\item L'utente seleziona il comando per caricare il file SVG.
\item L'utente seleziona il file SVG.
\item Il sistema visualizza il file SVG nel canvas. Se il file contiene elementi o attributi non supportati, essi saranno ignorati.
\end{enumerate}
\paragraph{Flusso alternativo}
Se il file non pu\`o essere caricato o a causa di un errore in lettura o perch\`e non \'e un file SVG, l'utente viene avvisato dell'errore e l'applicazione ritorna allo stato precedente l'esecuzione del comando.
\paragraph{Precondizioni} Il canvas \'e visualizzato.
\paragraph{Postcondizioni} Il canvas contiene la visualizzazione del file SVG selezionato dall'utente ed \'e modificabile.

\subsubsezione{UC9 Caso d'uso: disegna}
Si rimanda a \ref{ucdisegna} per una descrizione approfondita.

\subsubsezione{UC10 Caso d'uso: seleziona e modifica}
Si rimanda a \ref{ucselezionaemodifica} per una descrizione approfondita.

\subsezione{Vista UC8 - Disegna}
\label{ucdisegna}
Il seguente grafico illustra le interazioni tra l'utente e il sistema per disegnare un oggetto sul canvas.
\label{ucdisegna}

\subsubsezione{UC8.1 Caso d'uso: disegna linea retta/curva di Bezier}
\paragraph{Attori coinvolti} Utente.
\paragraph{Scopo e descrizione sintetica} 
Il sistema permette all'utente di disegnare una linea retta o una curva di Bezier.
\paragraph{Flusso di eventi}
\begin{elenconumerato}[\textbf{}]{\subsubsecindent}
\item L'utente seleziona il comando per disegnare una linea retta o una curva di Bezier.
\item L'utente disegna sul canvas la linea retta con due clic distinti, oppure la curva di Bezier tenendo premuto il pulsante sinistro del mouse.
\item Il sistema disegna la linea/curva.
\end{elenconumerato}
\paragraph{Precondizioni} Il canvas \'e visualizzato.
\paragraph{Postcondizioni} Il canvas mostra la nuova retta/curva disegnata.

\subsubsezione{UC8.2 Caso d'uso: disegna quadrato/rettangolo}
\paragraph{Attori coinvolti} Utente.
\paragraph{Scopo e descrizione sintetica} 
Il sistema permette all'utente di disegnare un quadrato o un rettangolo di qualsiasi dimensione.
\paragraph{Flusso di eventi}
\begin{elenconumerato}[\textbf{}]{\subsubsecindent}
\item L'utente seleziona il comando per disegnare un quadrato o un rettangolo.
\item L'utente disegna sul canvas il quadrato o il rettangolo, scegliendo col cursore la sua dimensione.
\item Il sistema disegna il quadrato/rettangolo.
\end{elenconumerato}
\paragraph{Precondizioni} Il canvas \'e visualizzato.
\paragraph{Postcondizioni} Il canvas mostra il nuovo quadrato/rettangolo disegnato.

\subsubsezione{UC8.3 Caso d'uso: disegna poligono}
\paragraph{Attori coinvolti} Utente.
\paragraph{Scopo e descrizione sintetica} 
Il sistema permette all'utente di disegnare un poligono di qualsiasi dimensione.
\paragraph{Flusso di eventi}
\begin{elenconumerato}[\textbf{}]{\subsubsecindent}
\item L'utente seleziona il comando per disegnare un quadrato o un rettangolo.
\item L'utente disegna sul canvas il quadrato o il rettangolo, scegliendo col cursore la sua dimensione.
\item Il sistema disegna il quadrato/rettangolo.
\end{elenconumerato}
\paragraph{Precondizioni} Il canvas \'e visualizzato.
\paragraph{Postcondizioni} Il canvas mostra il nuovo quadrato/rettangolo disegnato.

\subsubsezione{UC8.4 Caso d'uso: disegna casella di testo}
\paragraph{Attori coinvolti} Utente.
\paragraph{Scopo e descrizione sintetica} 
Il sistema permette all'utente di disegnare una casella di testo.
\paragraph{Flusso di eventi}
\begin{elenconumerato}[\textbf{}]{\subsubsecindent}
\item L'utente seleziona il comando per disegnare una casella di testo.
\item L'utente disegna sul canvas la casella, scegliendo col cursore la sua dimensione.
\item L'utente scrive il testo all'interno della casella.
\item Il sistema disegna la casella di testo.
\end{elenconumerato}
\paragraph{Precondizioni} Il canvas \'e visualizzato.
\paragraph{Postcondizioni} Il canvas mostra la nuova casella di testo disegnata.
%inizio requisiti
\sezione{Lista dei requisiti}
\label{listarequisiti}
I requisiti qui elencati sono suddivisi per tipologia (funzionale, prestazionale, di qualit\`a, di interfacciamento, d'ambiente); 
per ogni tipologia \`e presente un'ulteriore classificazione per classe di spazio di negoziato (obbligatorio, desiderabile, facoltativo). 
\subsezione{Requisiti funzionali (RF)}
\subsubsezione{Obbligatori (RFO)}
\begin{elenconumerato}[\textbf{RFO-}]{\subsubsecindent}
\item{AJAXDRAW permette di disegnare linee rette.}
\item{AJAXDRAW permette di disegnare curve.}
\item{AJAXDRAW permette di effettuare uno zoom in avanti e all'indietro nel disegno.}
\item{Per ogni elemento disegnato, AJAXDRAW permette all'utente di selezionare e spostare tale elemento.}
\item{AJAXDRAW permette di disegnare quadrilateri di qualsiasi dimensione(quadrata o rettangolare).}
\item{AJAXDRAW permette all'utente di disegnare cerchi o ellissi.}
\item{AJAXDRAW permette di disegnare poligoni regolari di qualsiasi dimensione e numero di lati.}
\item{AJAXDRAW permette di aggiungere caselle di testo al documento. }
\item{AJAXDRAW permette di cambiare il colore di riempimento, del bordo, delle linee, dell'opacità degli oggetti disegnati.}
\item{AJAXDRAW permette il ridimensionamento degli oggetti disegnati.}
\item{AJAXDRAW permette all'utente di selezionare ognuna delle funzionalit\`a suddette tramite una comoda barra degli strumenti.}
\item{Qualora suddette funzioni siano regolate da parametri, AJAXDRAW permette all'utente di impostarli.}
\item{AJAXDRAW \`e in grado di salvare le immagini create dall'utente nel formato SVG.}
\item{AJAXDRAW permette di caricare immagini in formato SVG che utilizzano il sottoinsieme che AJAXDRAW \'e in grado di generare.}
\end{elenconumerato}

%requisiti desiderabili
\subsubsezione{Desiderabili (RFD)}
\begin{elenconumerato}[\textbf{RFD-}]{\subsubsecindent}
\item{AJAXDRAW permette di disegnare a mano libera.}
\item{AJAXDRAW permette la creazione di spirali.}
\item{AJAXDRAW permette di utilizzare il disegno calligrafico o pennellato.}
\item{AJAXDRAW permette di disegnare connettori per diagrammi.}
\item{AJAXDRAW permette di disegnare curve di Bezier.}
\end{elenconumerato}

%facoltativi
\subsubsezione{Facoltativi (RFF)}
\begin{elenconumerato}[\textbf{RFF-}]{\subsubsecindent}
\end{elenconumerato}

%requisiti prestazionali
\subsezione{Requisiti prestazionali}
\subsubsezione{Obbligatori (RPO)}
\begin{elenconumerato}[\textbf{RPO-}]{\subsubsecindent}
\item L'utente non deve notare rallentamenti durante l'utilizzo dell'applicazione che non siano direttamente imputabili alla lentezza della connessione.
\end{elenconumerato}

%Requisiti di Qualità
\subsezione{Requisiti di qualit\`a (RQ)}	
\subsubsezione{Obbligatori (RQO)}
\begin{elenconumerato}[\textbf{RQO-}]{\subsubsecindent}
\item Il lato client dell'applicativo dev'essere realizzato usando tecnologie web standard\footnote{con il termine ``standard'' si intendono tecnologie o effettivamente standardizzate come l'\underline{ECMAScript} o tecnologie considerate standard \textitalic{de facto} come quelle raccomandata dal \textitalic{W3C}}.
\item Il codice sorgente rispetta lo standard di codifica interno all'azienda.
\item Il codice sorgente \`e scritto in accordo con i paradigmi dell'orientamento agli oggetti.
\item Il codice sorgente \`e corredato di documentazione completa.
\item La progettazione sar\`a orientata alla flessibilit\`a del codice.
\item Lo sviluppo del prodotto sar\`a accompagnato da fasi di test opportunamente documentate.
\item Il prodotto non necessita di nessuna installazione, basandosi su una pagina web HTML.
\item Il server dev'essere in grado di gestire pi\`u utenti in parallelo in modo indipendente.
\item L'eventuale richiesta di uscita dal programma da parte dell'utente deve avere sempre la precedenza sulle altre attivit\`a che il programma deve compiere.
\item Il funzionamento dell'applicazione dev'essere indipendente dalla risoluzione dello schermo dell'utente.
\end{elenconumerato}

%requisiti di qualita' desiderabili
%\subsubsezione{Desiderabili (RQD)}
%\begin{elenconumerato}[\textbf{RQD-}]{\subsubsecindent}
%\item 

%\end{elenconumerato}

%inizio requisiti di interfacciamento
\subsezione{Requisiti di interfacciamento (RI)}
\subsubsezione{Obbligatori (RIO)}
\begin{elenconumerato}[\textbf{RIO-}]{\subsubsecindent}
\item L'utilizzo del prodotto risulta semplice ed intuitivo grazie ad un'interfaccia grafica basata sull'uso di una toolbar.
\item La guida in linea \`e completa e di facile utilizzo per permettere all'utente di apprendere velocemente le funzionalit\`a del prodotto.
\item La guida in linea \`e richiamabile in qualsiasi momento.
\item L'uso della guida in linea non influenza il corretto funzionamento del programma. 
\end{elenconumerato}

%inizio requisiti d'ambiente
\subsezione{Requisiti d'ambiente (RA)}
\subsubsezione{Obbligatori (RAO)}
\begin{elenconumerato}[\textbf{RAO-}]{\subsubsecindent}
\item Il prodotto \`e pensato per funzionare su un internet browser.
\item Il prodotto dovr\`a essere implementato seguendo gli standard del linguaggio di markup HTML 5.
\item Il prodotto in particolare dovr\` a utilizzare il tag "canvas" appositamente creato per l'elaborazione del disegno vettoriale sul web.
\item Il software di supporto per lo sviluppo del prodotto \`e di tipo Open Source e freeware.
\item Il prodotto realizzato dovr\`a essere pubblicato sul sito "www.sourceforge.net", in conformit\`a con i relativi requisiti di natura Open Source, per favorire la continuit\`a del prodotto risultante.
\item Il prodotto \`e sviluppato e testato sui principali internet browsers: Mozilla Firefox, Google Chrome, Opera ed Apple Safari.
\end{elenconumerato}
\subsubsezione{Desiderabili (DAF)}
\begin{elenconumerato}[\textbf{DAF-}]{\subsubsecindent}
\item{AJAXDRAW garantisce la compatibilit\`a con il browser Microsoft Internet Explorer, il quale non implementa a pieno HTML 5, tramite la libreria \H{}excanvas\H{} creata da Google.}
\end{elenconumerato}


\newpage
\end{document}